% !TEX root = ../nc-cubic.tex
\subsection{二変数因子の同値と巡回五頂点型の可換実現}
\label{nqr:sec:roots}
本節では実際に圏同値を構成する。直前の次元検査は、先に構成した
同値における群作用の合成則を判定する箇所だけに用いる。

\begin{lemma}[二変数三次の巡回同変化]\label{nqr:lem:binary}
$\mathcal B=\MF^{\mathbb Z}(x^3+y^3)$ とし、$H=\mathbb Z/3$ の生成元は
$(x,y)\mapsto(\zeta x,\zeta^{-1}y)$ で作用するとする。
$\mathcal B^H\simeq\mathcal B$ という dg Morita 同値が存在し、
内部 shift と、両辺で指定した $(3)\simeq[2]$ の周期を保つように取れる。
\end{lemma}
\begin{proof}
$R=\C[x,y]/(x^3+y^3)$、三つの線形因子を $l_i$、$M_i=R/(l_i)$ とする。
$k=R/(x,y)$ と各 $M_i$ を次数付き特異圏で考える。
Orlov の parameter $-1$、window $i=0$ の分解
\cite[Theorem~2.5(ii) and its proof]{NFMOrlovSingularities}は
$k$ と三点の切断加群の次数 $\ge1$ の部分で全圏を生成する。
各点における非零一次式の乗法により、この部分は $M_i(-1)$ と同型である。
$0\to M_i(-1)\to M_i\to k\to0$ から $k,M_1,M_2,M_3$ も生成する。

$M_i$ の完全自由解消は $l_i,(x^3+y^3)/l_i$ が交互に作用し、
第 $2m,2m+1$ 項の生成次数は $3m,3m+1$ である。従って
$\RHom(M_i,k)=\C$ は次数零のみで、$\RHom(M_i,k(1))=0$。
$M_j$ に値を取る次数零 Hom 複体では、非零な線形または二次式は
$\C[t]$ の対応する二つの非負次数成分の間の同型となる。
$i=j$ のとき線形微分は零、$i\ne j$ のとき二次微分が零であり、
$\RHom(M_i,M_j)=\delta_{ij}\C$ を得る。
$k$ は上の分解で例外的である。
Serre 関手は $(1)$ なので $\RHom(k,M_i)=0$ でもある。
実際、\cite[Theorem~2.18]{NQRFaveroKelly} の同変 Serre 関手は
この二変数の場合 $(-2)[2]=(1)$ であり、$H$ を付けても
$dx\wedge dy$ の $H$-指標は自明で同じ式が成立する。
同定理の核群はここでは有限群、臨界軌跡は原点である。

以上より $\oplus M_i\oplus k$ は tilting 対象であり、
矢は三つの自然な剰余写像 $M_i\to k$ である。
$H$ は三つの枝を巡回し、$k$ には標準線型化がある。
$\mathcal B^H$ では誘導対象 $U=\operatorname{Ind}_H M_1$ と
$k$ の三つの指標付き対象 $k_\chi$ が生成する。
これは忘却後の生成と、平均化による誘導忘却の直和因子表示から従う。
随伴で計算すると矢は $U\to k_\chi$ の三本、各一次元で、
自己恒等以外の他の Hom と非零次数の Hom は消える。
従って同変化によって星型四頂点 quiver の矢が逆転する。

元の圏でも $N=\operatorname{Cone}(\oplus M_i\to k)[-1]$ を使えば
同じ向きの tilting 表示を得る。三角から $\RHom(N,M_i)=\C$、
$\RHom(M_i,N)=0$、$\RHom(N,N)=\C$ である。
いずれの生成代数も次数零にのみコホモロジーを持つため
通常の切断によって道代数に dg 準同型の列で結ばれる。
$U\leftrightarrow N$, $k_\chi\leftrightarrow M_i$ が所要の同値を与える。
両辺で内部 shift が Serre 関手であることから、同値はその可逆双加群を
共役する。比較同型の三乗と指定周期との差は星型道代数の中心の単元
$\C^\times$ であり、三乗根による補正で周期を同期できる。
\end{proof}

\begin{corollary}[任意次元の分離した二変数対]\label{nqr:cor:separated-pairs}
$n\ge3$、$n+2=2t+m$, $t,m\ge1$ とし、$g$ を $m$ 変数の三次斉次式で
アフィン臨界点が原点だけのものとする。
$X=V(\sum_{a=1}^t(x_a^3+y_a^3)+g)\subset\PP^{n+1}$ と置く。
$H=(\mathbb Z/3)^t$ の第 $a$ 生成元は第 $a$ 対だけを
$(\zeta x_a,\zeta^{-1}y_a)$ に移す。このとき
$\Ku(X)^H\simeq\Ku(X)$ という $\C$-線型 dg Morita 同値がある。
\end{corollary}
\begin{proof}
分離和の勾配より $X$ は滑らかであり、Orlov の同値で全残余圏を
通常の整数次数付き MF で計算できる。切断、次数による切り詰めおよび
自由加群の window は各座標対の線型作用と可換する。
第零ブロックを $g$、残りを二変数対として
$\Lambda=\mathbb Z^{t+1}/\langle 3(e_a-e_0)\rangle_{a=1}^t$ と置く。
各ブロックの整数次数付き MF の Morita tensor 積は、分離和の
$\Lambda$ 次数付き MF である
\cite[Corollary~2.40 and (3.5)]{MBRDBFKFurther}。
外積の構成は各ブロックの引戻しと可換する。
$\Lambda\to\mathbb Z$ の核は $e_a-e_0$ の張る $(\mathbb Z/3)^t$ であり、
各生成元は二つの内部 shift の比として作用する。
この次数差の群で同変化すると整数次数付き MF が復元される。
実際、次数を粗くするのは次数差の shift による軌道圏の冪等完備化であり
\cite[Proposition~4.5, Corollary~4.6]{MBRDBFKFurther}、
有限群の位数が可逆なので誘導対象の Hom と平均化からこの圏は同変圏と同値である。
各比の三乗の自明化は両因子での $(3)=[2]$ の差で与えられる。
したがって同変化は各二変数因子での $H_a$ の同変化と交換できる。
$H_a$ は他の因子を固定するので、$H$-同変化した積は
$\mathcal B^{H_a}$ の各因子と $g$ の因子との積となる
（誘導生成対象の自己準同型でも従う）。
補題\ref{nqr:lem:binary}の同値を各因子で tensor する。
周期の整合性を具体化する。二変数の星型道代数の中心は
$\C$ であり、比較同型 $\gamma_a$ の三回合成が指定周期とずれる
比を $c_a\in\C^\times$ とすれば $c_a^{-1}$ の三乗根で
$\gamma_a$ を掛けて比を一にできる。次数差 $e_a-e_0$ の比較は
この同型と第零因子の恒等比較との外積である。二つの異なる次数差で
共有されるのは第零因子の同じ逆 shift の冪だけであり、他因子の
次数零の比較射の外積は交換する。従って各生成元の三乗関係と
全ての交換関係が保たれる。星型道代数の負次数 Hochschild
コホモロジーは零で、因子の比較同型間に高次の選択は残らない。
このように
すべての次数差の合成則を保つ。よって通常の次数に降下して結論を得る。
\end{proof}

\begin{lemma}[一つの差ベクトルでの群作用付き同値]
\label{nqr:lem:one-root}
$P,Y,\mathcal C$ を前節のものとする。$r=e_i-e_j$, $H=\langle r\rangle$、
$J=H^\perp\subset P$ とすれば $r\cdot r=2$ より $P=H\oplus J$ である。
$\mathcal C^H$ に $P$ を作用させる際、$a=a_H+a_J$ に対して
$a_J$ は幾何的作用、$a_H$ は $h\mapsto\zeta^{-a_H\cdot h}$ の指標作用を用いる。
この作用を持つ $\mathcal C^H$ と標準対角作用を持つ $\mathcal C$ の間に、
合成則まで保つ $P$-同変 dg Morita 同値が存在する。
\end{lemma}
\begin{proof}
$r=e_0-e_1$ としてよい。系\ref{nqr:cor:separated-pairs}に
$t=1$, $g=x_2^3+x_3^3+x_4^3$ を代入し、
群作用を忘れた同値 $\mathcal C^H\simeq\mathcal C$ を得る。

次に $P$ を運ぶ。$HH_1(\mathcal C^H)$ の恒等 sector から
$b_2,b_3,b_4$ の三つの指標が来る。$u=\pm r$ の各 sector は残り三座標の
平面楕円曲線であり、その指標指数は
$\sum_{j=2}^4(b_J)_j-b_H\cdot u=(b_J)_0-b_H\cdot u$ だから、
残る二つは $b_0,b_1$ である。
よって忠実な $P$-作用である。Liu の定理と
命題\ref{nqt:prop:fermat-target}の同時対角化により、輸送像を
Fermat 三様体の標準対角群へ共役できる。
対角座標の再尺度化で Fermat 方程式を戻せるので、共役は両端の
Fermat 三様体の幾何的同型で行える。五つの $HH_1$ 指標が標準のものに
一致するため、座標置換の後に、個々の $b$ の作用は標準の $b$ と同型になる。
核の一意性とスカラーの合成障害は命題\ref{nqs:prop:plane-coherence}の
証明と同じである。

最後に、こうして得た作用の合成のずれ $B_H$ が零であることを確認する。
前節と同じ固定点公式で、今回は
\[
 F_i^H(a,b)=\frac13\sum_{u,v\in H}\zeta^{a\cdot v-b\cdot u}
 R_i(a_J+u,b_J+v)
\]
を用いる。二段目の部分群 $K$ に対する平均公式は
$H\times K$ が $Y$ に $(u,a)\mapsto u+a_J$ を通して作用し、
交換子に $a\cdot v-b\cdot u$ が付くことから従う。
ここで $F_i^H=R_i$ を全体の表の列挙を使わずに示す。
$z$ を形式変数、$x,v\in\mathbb F_3$ とし
\[
 K(x,v)=
 \begin{cases}
 z\zeta^v+z^{-1}\zeta^{-v}&x=0,\\
 1&x\ne0
 \end{cases}
\]
と置くと、前節の固定点公式は
\[
 R_i(g,h)=[z^{3i}]\sum_{d\in\mathbb F_3}\prod_{j=0}^4K(g_j-d,h_j)
\]
と書ける。実際、空でないレベル集合 $I$ では
$e_j\in\{1,2\}$ を選ぶ原始留数項の指標は $\zeta^{\sum e_j h_j}$、
Hodge 次数差は $|I|-2\sum e_j/3$ である。これは Fermat 勾配商の
単項式 $\prod x_j^{e_j-1}$ の留数で、整数次数への係数抽出は
$\sum e_j\equiv0\pmod3$ を課す。原始部分以外の $HH_0$ は
非空な各レベルから $|I|-1$ 次元あり、例外二項の減算後は
$5-\#\{I\ne\varnothing\}-2$、すなわち空のレベルの個数になる。
これが空積の $1$ に一致する。

$d,a_J,b_J$ を固定し $A=(a_J)_0-d$, $B=(b_J)_0$, $E=z\zeta^B$
と書く。$u=kr$, $v=lr$, $a_H=sr$, $b_H=tr$ とすると、
動く二座標の積は
$D(k,l)=K(A+k,B+l)K(A-k,B-l)$ である。残る三因子は $k,l$ に依らない。
$r\cdot r=2=-1$ より有限 Fourier 変換の核は
$\zeta^{-sl+tk}/3$ であり、
$\mathbf1_{k=c}\zeta^{ql}$ を $\mathbf1_{s=q}\zeta^{ct}$ に送る。
よってこの基底で係数が $c,q$ に対称であればよい。

$A\ne0$ のとき $c=q=0$ の係数は $1$、
$c,q\ne0$ の係数は $E^m$（$m\in\{1,-1\}$ は $m=-Acq\bmod3$）で、
他の係数は零である。実際、$k=c\ne0$ では一座標だけが寄与し、
指数は $B-Acl$ である。$A=0$ のときは
\[
 D(k,l)=1+(E^2+E^{-2}-2)\mathbf1_{k=0}
          +3\mathbf1_{k=l=0}.
\]
この変換は定数 $1$ と $3\mathbf1_{k=l=0}$ を交換し、$\mathbf1_{k=0}$
を固定する。従って両場合とも変換後は $D(s,t)$ である。
不動の三因子を掛けて $d$ の和と係数抽出を行えば所要の恒等式が従う。
合成障害を検出するには六平面の $HH_1$ だけで足りる。
$P$ の基底 $e_i-e_4$ ($0\le i<4$) の各二元が張る平面を $L$ とする。
三座標だけが動き、それらの三指標は非零で和が零、他の二座標は固定される。
まず恒等 sector の不変一次留数は二次元である。
非零元のうち、三座標の指数がすべて等しい二元では、その三座標の楕円曲線の
留数指標が自明である。従って乗数なしでは各一次元を与え、非零交代形式で
ねじると両方消える。他の六元は三対の逆元に分かれる。各対は固定二座標と
指数零の一座標からなる楕円曲線を固定し、留数はその一座標の非自明指標である。
この指標の核は当該の位数三部分群である。非零交代形式の交換子指標は、
対の二元でこの核を持つ二つの非自明指標を一度ずつ取る。
従って乗数なしでは寄与は零、非零のねじれでは各対から一だけ残る。
合計は順に $2+2=4$ と $2+3=5$ である。
先の $F_1^H=R_1$ の恒等式を $L$ で平均すると輸送側のねじれなしの値も四である。
従って $B_H$ は六平面すべてに零として制限される。
六つの基底対の値が交代形式を決定するので $B_H=0$。
これは $130$ 平面・$729$ 形式での総当たりとも一致する。
これは既に存在する同値についての合成障害の計算である。
$B_H=0$ なら相違するスカラー cocycle は coboundary なので、
比較同型をその $1$-cochain で補正すれば群合成を保つ。
負の Hochschild コホモロジーが消えているため、核の同型の間に
さらに選択すべき高次ホモトピーはない。
\end{proof}

\begin{lemma}[差ベクトルを加えるときの乗数]\label{nqr:lem:enlargement}
$K\subset P$, $H\cap K=0$ とし、$B_K$ を $K$ 上の乗数の交換子形式とする。
$\mathcal C^{K,B_K}$ は $\mathcal C^{H+K,B'}$ と同値である。
ここでは上付き形式はその形式を持つ任意の乗数を表し、
$B'|_K=B_K$, $B'(r,k)=-r\cdot k$ とする。
\end{lemma}
\begin{proof}
$K$ の $J$ への射影を $S$ とし、$K=\{s+\ell(s):s\in S\}$ と書く。
$H+K=H\oplus S$ 上で、$S$ に $B_K$ を移した形式と、
$B'(a,s)=-a\cdot\ell(s)$ ($a\in H$) を使う。
$K$ に戻すと交差項が内積の対称性で相殺され、表示の $B'$ となる。

$(\mathcal C^H)^{H\oplus S,B'}$ を二段の線型化として書く。
第一段の $U_h$ は幾何的 $H$-作用、第二段の $V_a$ は指標作用の線型化、
$W_s$ は幾何的 $S$-作用の線型化を表す。
$B'$ の代表は $H$ 上自明に選び、$V_a$ を通常の可換な位数三の作用素とする。
$V_aU_h=\zeta^{a\cdot h}U_hV_a$ および
$V_aW_s=\zeta^{-a\cdot\ell(s)}W_sV_a$ である
（幾何的作用の pullback を含めた合成の記法）。
$W_s$ は第一段の $H$-線型化と可換する。
$V_a$ の同時固有成分を $d\in H$、固有指標を $a\mapsto\zeta^{a\cdot d}$
で表示すると、$U_h$ は成分を $h$ だけ、$W_s$ は $-\ell(s)$ だけ移す。
そこで成分 $0$ には $U_{\ell(s)}W_s$ が作用する。
幾何的作用は $s+\ell(s)$、その積の乗数は指定した $K$ の乗数である。
この成分を取る操作は同値である。逆は成分 $0$ から全 $H$-移動の
直和を誘導することであり、線型化は上記交換則と $K$-線型化で与えられる。
これは有限群の推移的な成分集合に対する安定化部分群との誘導・制限で、
標数零の指標冪等元は dg Morita 圏でも分裂する。
従って二段圏は $\mathcal C^{K,B_K}$ である。
一方、補題\ref{nqr:lem:one-root}により同じ二段圏は
$\mathcal C^{H\oplus S,B'}$ に同値となる。
\end{proof}

\begin{theorem}[巡回五頂点型の Fermat 実現]\label{nqr:thm:cyclic-fermat}
巡回五頂点型の $\mathcal Q$ は可換 Fermat 三様体の $\mathcal C$ と
$\C$-線型 dg Morita 同値である。特に定理\ref{nqt:thm:reduction}の
五変数・交換係数が三乗根の全表示族で狭い主予想が成立する。
\end{theorem}
\begin{proof}
次の順で差ベクトルを取る。
\[
 r_1=e_0-e_1,\quad r_2=e_0-e_3,\quad
 r_3=e_1-e_2,\quad r_4=e_1-e_4 .
\]
五頂点を結ぶ木の四辺なので $P$ の基底である。
零部分群から補題\ref{nqr:lem:enlargement}を四回適用すると
最終形式は $i<j$ で $B'(r_i,r_j)=r_i\cdot r_j$ となる。
この基底で $A(g,h)=(D^{-1}g)\cdot h$ と直接比較すると、両者はともに
\[
 \begin{pmatrix}0&1&-1&-1\\-1&0&0&0\\1&0&0&1\\1&0&-1&0\end{pmatrix}
 \pmod3
\]
である。よって四つの同値の合成は
$\mathcal C\simeq\mathcal C^{P,\alpha}=\mathcal Q$ を与える。
定理\ref{nqt:thm:reduction}と合わせれば全表示族で結論する。
他の周囲代数や他次元の全対象への拡張は未解決である。
\end{proof}
